V ďalšom kroku bol stiahnutý a nainštalovaný nástroj Dorado. Na základe vlastností datasetu pod5 bol zvolený vhodný model Dorado určený na spracovanie daného datasetu.
Výber modelu vychádzal z typu nanopórového zariadenia, pomocou ktorého bol dataset získaný.
Po spracovaní basecallerom Dorado vznikol BAM súbor obsahujúci preložené sekvencie.
Spracovanie trvalo približne X hodín, v závislosti od veľkosti datasetu a výpočtovej kapacity použitého hardvéru. %
